% =========================================================================
% REFERENCIAS
% =========================================================================
% REFERENCIAS BIBLIOGRÁFICAS (APA 7ª edición)
\addcontentsline{toc}{chapter}{Referencias Bibliográficas}

\begin{thebibliography}{99}
\bibitem{akiba2019} Akiba, T., Sano, S., Yanase, T., Ohta, T., \& Koyama, M. (2019). Optuna: A next-generation hyperparameter optimization framework. In \textit{Proceedings of the 25th ACM SIGKDD International Conference on Knowledge Discovery \& Data Mining} (pp. 2623--2631). ACM. \url{https://doi.org/10.1145/3292500.3330701}
\bibitem{altman1968} Altman, E. I. (1968). Financial ratios, discriminant analysis and the prediction of corporate bankruptcy. \textit{The Journal of Finance, 23}(4), 589--609. \url{https://doi.org/10.1111/j.1540-6261.1968.tb00843.x}
\bibitem{arlot2010} Arlot, S., \& Celisse, A. (2010). A survey of cross-validation procedures for model selection. \textit{Statistics Surveys, 4}, 40--79. \url{https://doi.org/10.1214/09-SS054}
\bibitem{armendariz2010} Armendáriz, B., \& Morduch, J. (2010). \textit{The economics of microfinance} (2nd ed.). MIT Press.
\bibitem{banxico2025} Banco de México. (2025). Determinantes de la morosidad de la cartera de consumo para las sofipos [Recuadro 1]. En \textit{Reporte de Estabilidad Financiera, diciembre de 2025} (pp. 33--36). \url{https://www.banxico.org.mx/publicaciones-y-prensa/reportes-sobre-el-sistema-financiero/recuadros/%7BBF224D96-2D2C-22A1-328E-2D98D71E3532%7D.pdf}
\bibitem{banxicosie} Banco de México. (s. f.). \textit{Sistema de Información Económica (SIE)}. \url{https://www.banxico.org.mx/SieInternet/}% VERIFICAR: fecha de consulta y series exactas (TIIE a 28 días, Cetes a 28 días, tipo de cambio FIX) en el SIE.
\bibitem{barboza2017} Barboza, F., Kimura, H., \& Altman, E. (2017). Machine learning models and bankruptcy prediction. \textit{Expert Systems with Applications, 83}, 405--417. \url{https://doi.org/10.1016/j.eswa.2017.04.006}
\bibitem{basel2011} Basel Committee on Banking Supervision. (2011). \textit{Basel III: A global regulatory framework for more resilient banks and banking systems} (rev. ed.). Bank for International Settlements. \url{https://www.bis.org/publ/bcbs189.htm}
\bibitem{batista2004} Batista, G. E. A. P. A., Prati, R. C., \& Monard, M. C. (2004). A study of the behavior of several methods for balancing machine learning training data. \textit{ACM SIGKDD Explorations Newsletter, 6}(1), 20--29. \url{https://doi.org/10.1145/1007730.1007735}
\bibitem{bellman1957} Bellman, R. (1957). \textit{Dynamic programming}. Princeton University Press.% VERIFICAR: datos de edición del libro (no registrado en CrossRef).
\bibitem{bengio1994} Bengio, Y., Simard, P., \& Frasconi, P. (1994). Learning long-term dependencies with gradient descent is difficult. \textit{IEEE Transactions on Neural Networks, 5}(2), 157--166. \url{https://doi.org/10.1109/72.279181}
\bibitem{bergmeir2018} Bergmeir, C., Hyndman, R. J., \& Koo, B. (2018). A note on the validity of cross-validation for evaluating autoregressive time series prediction. \textit{Computational Statistics \& Data Analysis, 120}, 70--83. \url{https://doi.org/10.1016/j.csda.2017.11.003}
\bibitem{bergstra2012} Bergstra, J., \& Bengio, Y. (2012). Random search for hyper-parameter optimization. \textit{Journal of Machine Learning Research, 13}, 281--305. \url{https://jmlr.org/papers/v13/bergstra12a.html}
\bibitem{beutel2019} Beutel, J., List, S., \& von Schweinitz, G. (2019). Does machine learning help us predict banking crises? \textit{Journal of Financial Stability, 45}, Article 100693. \url{https://doi.org/10.1016/j.jfs.2019.100693}
\bibitem{bottou2018} Bottou, L., Curtis, F. E., \& Nocedal, J. (2018). Optimization methods for large-scale machine learning. \textit{SIAM Review, 60}(2), 223--311. \url{https://doi.org/10.1137/16M1080173}
\bibitem{box1976} Box, G. E. P., \& Jenkins, G. M. (1976). \textit{Time series analysis: Forecasting and control} (rev. ed.). Holden-Day.% VERIFICAR: datos de edición del libro (no registrado en CrossRef).
\bibitem{boyd2004} Boyd, S., \& Vandenberghe, L. (2004). \textit{Convex optimization}. Cambridge University Press. \url{https://doi.org/10.1017/CBO9780511804441}
\bibitem{breiman1996} Breiman, L. (1996). Bagging predictors. \textit{Machine Learning, 24}(2), 123--140. \url{https://doi.org/10.1007/BF00058655}
\bibitem{breiman2001} Breiman, L. (2001a). Random forests. \textit{Machine Learning, 45}(1), 5--32. \url{https://doi.org/10.1023/A:1010933404324}
\bibitem{breiman2001cultures} Breiman, L. (2001b). Statistical modeling: The two cultures. \textit{Statistical Science, 16}(3), 199--231. \url{https://doi.org/10.1214/ss/1009213726}
\bibitem{breiman1984} Breiman, L., Friedman, J. H., Olshen, R. A., \& Stone, C. J. (1984). \textit{Classification and regression trees}. Wadsworth.% VERIFICAR: datos de edición del libro (no registrado en CrossRef).
\bibitem{brier1950} Brier, G. W. (1950). Verification of forecasts expressed in terms of probability. \textit{Monthly Weather Review, 78}(1), 1--3. \url{https://doi.org/10.1175/1520-0493(1950)078<0001:VOFEIT>2.0.CO;2}
\bibitem{bussmann2021} Bussmann, N., Giudici, P., Marinelli, D., \& Papenbrock, J. (2021). Explainable machine learning in credit risk management. \textit{Computational Economics, 57}(1), 203--216. \url{https://doi.org/10.1007/s10614-020-10042-0}
\bibitem{canomorales2023} Cano Morales, A. M., y Torres Grimaldo, J. A. (2023). El gobierno corporativo y sus efectos en el índice de solvencia de las sociedades financieras populares en México. \textit{Brazilian Journal of Business, 5}(1), 271--289. \url{https://doi.org/10.34140/bjbv5n1-017}
\bibitem{casella2002} Casella, G., \& Berger, R. L. (2002). \textit{Statistical inference} (2nd ed.). Duxbury.% VERIFICAR: datos de edición del libro (no registrado en CrossRef).
\bibitem{cerqueira2020} Cerqueira, V., Torgo, L., \& Mozetič, I. (2020). Evaluating time series forecasting models: An empirical study on performance estimation methods. \textit{Machine Learning, 109}(11), 1997--2028. \url{https://doi.org/10.1007/s10994-020-05910-7}
\bibitem{chawla2002} Chawla, N. V., Bowyer, K. W., Hall, L. O., \& Kegelmeyer, W. P. (2002). SMOTE: Synthetic minority over-sampling technique. \textit{Journal of Artificial Intelligence Research, 16}, 321--357. \url{https://doi.org/10.1613/jair.953}
\bibitem{chen2016} Chen, T., \& Guestrin, C. (2016). XGBoost: A scalable tree boosting system. In \textit{Proceedings of the 22nd ACM SIGKDD International Conference on Knowledge Discovery and Data Mining} (pp. 785--794). \url{https://doi.org/10.1145/2939672.2939785}
\bibitem{chernozhukov2018} Chernozhukov, V., Chetverikov, D., Demirer, M., Duflo, E., Hansen, C., Newey, W., \& Robins, J. (2018). Double/debiased machine learning for treatment and structural parameters. \textit{The Econometrics Journal, 21}(1), C1--C68. \url{https://doi.org/10.1111/ectj.12097}
\bibitem{chicco2020} Chicco, D., \& Jurman, G. (2020). The advantages of the Matthews correlation coefficient (MCC) over F1 score and accuracy in binary classification evaluation. \textit{BMC Genomics, 21}, Article 6. \url{https://doi.org/10.1186/s12864-019-6413-7}
\bibitem{chicco2021} Chicco, D., Warrens, M. J., \& Jurman, G. (2021). The Matthews correlation coefficient (MCC) is more informative than Cohen's kappa and Brier score in binary classification assessment. \textit{IEEE Access, 9}, 78368--78381. \url{https://doi.org/10.1109/ACCESS.2021.3084050}
\bibitem{cnbv} Comisión Nacional Bancaria y de Valores. (s. f.). \textit{Boletín estadístico de Sociedades Financieras Populares}. Gobierno de México.
\bibitem{cinif} Consejo Mexicano de Normas de Información Financiera. (s. f.). \textit{NIF B-3, Estado de resultado integral}. CINIF. \url{https://www.cinif.org.mx/}% VERIFICAR: edición vigente de la NIF B-3 y su aplicación a las SOFIPOs, que se rigen por los criterios contables de la CNBV.
\bibitem{cormen2022} Cormen, T. H., Leiserson, C. E., Rivest, R. L., \& Stein, C. (2022). \textit{Introduction to algorithms} (4th ed.). MIT Press.% VERIFICAR: datos de edición del libro (no registrado en CrossRef).
\bibitem{davis2006} Davis, J., \& Goadrich, M. (2006). The relationship between precision-recall and ROC curves. In \textit{Proceedings of the 23rd International Conference on Machine Learning} (pp. 233--240). ACM. \url{https://doi.org/10.1145/1143844.1143874}
\bibitem{diaz2026} Díaz Mondragón, M. (2026). Sociedades financieras populares e inteligencia artificial. \textit{Revista Collage, 1}(1). \url{https://doi.org/10.67470/gi7713-cv6crv6f-mc}% VERIFICAR: páginas o número de artículo no registrados en CrossRef.
\bibitem{dickey1979} Dickey, D. A., \& Fuller, W. A. (1979). Distribution of the estimators for autoregressive time series with a unit root. \textit{Journal of the American Statistical Association, 74}(366a), 427--431. \url{https://doi.org/10.1080/01621459.1979.10482531}
\bibitem{dumitrescu2022} Dumitrescu, E., Hué, S., Hurlin, C., \& Tokpavi, S. (2022). Machine learning for credit scoring: Improving logistic regression with non-linear decision-tree effects. \textit{European Journal of Operational Research, 297}(3), 1178--1192. \url{https://doi.org/10.1016/j.ejor.2021.06.053}
\bibitem{fawcett2006} Fawcett, T. (2006). An introduction to ROC analysis. \textit{Pattern Recognition Letters, 27}(8), 861--874. \url{https://doi.org/10.1016/j.patrec.2005.10.010}
\bibitem{fernandez2018} Fernández, A., García, S., Herrera, F., \& Chawla, N. V. (2018). SMOTE for learning from imbalanced data: Progress and challenges, marking the 15-year anniversary. \textit{Journal of Artificial Intelligence Research, 61}, 863--905. \url{https://doi.org/10.1613/jair.1.11192}
\bibitem{fisher1930} Fisher, I. (1930). \textit{The theory of interest}. Macmillan.% VERIFICAR: datos de edición del libro (no registrado en CrossRef).
\bibitem{friedman2001} Friedman, J. H. (2001). Greedy function approximation: A gradient boosting machine. \textit{The Annals of Statistics, 29}(5), 1189--1232. \url{https://doi.org/10.1214/aos/1013203451}
\bibitem{fuster2022} Fuster, A., Goldsmith-Pinkham, P., Ramadorai, T., \& Walther, A. (2022). Predictably unequal? The effects of machine learning on credit markets. \textit{The Journal of Finance, 77}(1), 5--47. \url{https://doi.org/10.1111/jofi.13090}
\bibitem{geman1992} Geman, S., Bienenstock, E., \& Doursat, R. (1992). Neural networks and the bias/variance dilemma. \textit{Neural Computation, 4}(1), 1--58. \url{https://doi.org/10.1162/neco.1992.4.1.1}
\bibitem{gers2000} Gers, F. A., Schmidhuber, J., \& Cummins, F. (2000). Learning to forget: Continual prediction with LSTM. \textit{Neural Computation, 12}(10), 2451--2471. \url{https://doi.org/10.1162/089976600300015015}
\bibitem{grinsztajn2022} Grinsztajn, L., Oyallon, E., \& Varoquaux, G. (2022). Why do tree-based models still outperform deep learning on tabular data? \textit{arXiv}. \url{https://doi.org/10.48550/arXiv.2207.08815}
\bibitem{gu2020} Gu, S., Kelly, B., \& Xiu, D. (2020). Empirical asset pricing via machine learning. \textit{The Review of Financial Studies, 33}(5), 2223--2273. \url{https://doi.org/10.1093/rfs/hhaa009}
\bibitem{guo2017} Guo, C., Pleiss, G., Sun, Y., \& Weinberger, K. Q. (2017). On calibration of modern neural networks. \textit{arXiv}. \url{https://doi.org/10.48550/arXiv.1706.04599}
\bibitem{hamilton1994} Hamilton, J. D. (1994). \textit{Time series analysis}. Princeton University Press.% VERIFICAR: datos de edición del libro (no registrado en CrossRef).
\bibitem{harris2020} Harris, C. R., Millman, K. J., van der Walt, S. J., Gommers, R., Virtanen, P., Cournapeau, D., Wieser, E., Taylor, J., Berg, S., Smith, N. J., Kern, R., Picus, M., Hoyer, S., van Kerkwijk, M. H., Brett, M., Haldane, A., del Río, J. F., Wiebe, M., Peterson, P., \ldots{} Oliphant, T. E. (2020). Array programming with NumPy. \textit{Nature, 585}(7825), 357--362. \url{https://doi.org/10.1038/s41586-020-2649-2}
\bibitem{hastie2009} Hastie, T., Tibshirani, R., \& Friedman, J. (2009). \textit{The elements of statistical learning: Data mining, inference, and prediction} (2nd ed.). Springer. \url{https://doi.org/10.1007/978-0-387-84858-7}
\bibitem{he2009} He, H., \& Garcia, E. A. (2009). Learning from imbalanced data. \textit{IEEE Transactions on Knowledge and Data Engineering, 21}(9), 1263--1284. \url{https://doi.org/10.1109/TKDE.2008.239}
\bibitem{hochreiter1997} Hochreiter, S., \& Schmidhuber, J. (1997). Long short-term memory. \textit{Neural Computation, 9}(8), 1735--1780. \url{https://doi.org/10.1162/neco.1997.9.8.1735}
\bibitem{hoerl1970} Hoerl, A. E., \& Kennard, R. W. (1970). Ridge regression: Biased estimation for nonorthogonal problems. \textit{Technometrics, 12}(1), 55--67. \url{https://doi.org/10.1080/00401706.1970.10488634}
\bibitem{holopainen2017} Holopainen, M., \& Sarlin, P. (2017). Toward robust early-warning models: A horse race, ensembles and model uncertainty. \textit{Quantitative Finance, 17}(12), 1933--1963. \url{https://doi.org/10.1080/14697688.2017.1357972}
\bibitem{hosmer2013} Hosmer, D. W., Lemeshow, S., \& Sturdivant, R. X. (2013). \textit{Applied logistic regression} (3rd ed.). Wiley. \url{https://doi.org/10.1002/9781118548387}
\bibitem{hunter2007} Hunter, J. D. (2007). Matplotlib: A 2D graphics environment. \textit{Computing in Science \& Engineering, 9}(3), 90--95. \url{https://doi.org/10.1109/MCSE.2007.55}
\bibitem{hyndman2006} Hyndman, R. J., \& Koehler, A. B. (2006). Another look at measures of forecast accuracy. \textit{International Journal of Forecasting, 22}(4), 679--688. \url{https://doi.org/10.1016/j.ijforecast.2006.03.001}
\bibitem{inegiigae} Instituto Nacional de Estadística y Geografía. (s. f.-a). \textit{Indicador Global de la Actividad Económica (IGAE)}. INEGI. \url{https://www.inegi.org.mx/temas/igae/}% VERIFICAR: año base vigente y fecha de consulta.
\bibitem{inegiinpc} Instituto Nacional de Estadística y Geografía. (s. f.-b). \textit{Índice Nacional de Precios al Consumidor (INPC)}. INEGI. \url{https://www.inegi.org.mx/temas/inpc/}% VERIFICAR: año base vigente y fecha de consulta.
\bibitem{jordan2015} Jordan, M. I., \& Mitchell, T. M. (2015). Machine learning: Trends, perspectives, and prospects. \textit{Science, 349}(6245), 255--260. \url{https://doi.org/10.1126/science.aaa8415}
\bibitem{kaufman2012} Kaufman, S., Rosset, S., Perlich, C., \& Stitelman, O. (2012). Leakage in data mining. \textit{ACM Transactions on Knowledge Discovery from Data, 6}(4), 1--21. \url{https://doi.org/10.1145/2382577.2382579}% VERIFICAR: subtítulo del artículo no registrado en CrossRef.
\bibitem{ke2017} Ke, G., Meng, Q., Finley, T., Wang, T., Chen, W., Ma, W., Ye, Q., \& Liu, T.-Y. (2017). LightGBM: A highly efficient gradient boosting decision tree. In \textit{Advances in Neural Information Processing Systems 30} (pp. 3146--3154). \url{https://papers.nips.cc/paper/6907-lightgbm-a-highly-efficient-gradient-boosting-decision-tree}% VERIFICAR: autores, páginas y URL (actas de NeurIPS, sin DOI).
\bibitem{khandani2010} Khandani, A. E., Kim, A. J., \& Lo, A. W. (2010). Consumer credit-risk models via machine-learning algorithms. \textit{Journal of Banking \& Finance, 34}(11), 2767--2787. \url{https://doi.org/10.1016/j.jbankfin.2010.06.001}
\bibitem{kingma2015} Kingma, D. P., \& Ba, J. (2015). Adam: A method for stochastic optimization. \textit{arXiv}. \url{https://doi.org/10.48550/arXiv.1412.6980}
\bibitem{lajornada2017} \textit{La Jornada}. (2017, 30 de junio). Se dispara el índice de morosidad en las Sofipo. \url{https://www.jornada.com.mx/2017/06/30/economia/018n1eco}
\bibitem{lecun2015} LeCun, Y., Bengio, Y., \& Hinton, G. (2015). Deep learning. \textit{Nature, 521}(7553), 436--444. \url{https://doi.org/10.1038/nature14539}
\bibitem{lemaitre2017} Lemaître, G., Nogueira, F., \& Aridas, C. K. (2017). Imbalanced-learn: A Python toolbox to tackle the curse of imbalanced datasets in machine learning. \textit{Journal of Machine Learning Research, 18}(17), 1--5. \url{https://doi.org/10.48550/arXiv.1609.06570}
\bibitem{lessmann2015} Lessmann, S., Baesens, B., Seow, H.-V., \& Thomas, L. C. (2015). Benchmarking state-of-the-art classification algorithms for credit scoring: An update of research. \textit{European Journal of Operational Research, 247}(1), 124--136. \url{https://doi.org/10.1016/j.ejor.2015.05.030}
\bibitem{ley2001} Ley de Ahorro y Crédito Popular. (2001). Diario Oficial de la Federación. \url{https://www.diputados.gob.mx/LeyesBiblio/ref/lacp.htm}
\bibitem{lundberg2017} Lundberg, S. M., \& Lee, S.-I. (2017). A unified approach to interpreting model predictions. In \textit{Advances in Neural Information Processing Systems 30} (pp. 4765--4774). \url{https://papers.nips.cc/paper/7062-a-unified-approach-to-interpreting-model-predictions}
\bibitem{lundberg2020} Lundberg, S. M., Erion, G., Chen, H., DeGrave, A., Prutkin, J. M., Nair, B., Katz, R., Himmelfarb, J., Bansal, N., \& Lee, S.-I. (2020). From local explanations to global understanding with explainable AI for trees. \textit{Nature Machine Intelligence, 2}(1), 56--67. \url{https://doi.org/10.1038/s42256-019-0138-9}
\bibitem{mankiw2019} Mankiw, N. G. (2019). \textit{Macroeconomics} (10th ed.). Worth Publishers.% VERIFICAR: datos de edición del libro (no registrado en CrossRef).
\bibitem{matthews1975} Matthews, B. W. (1975). Comparison of the predicted and observed secondary structure of T4 phage lysozyme. \textit{Biochimica et Biophysica Acta (BBA), 405}(2), 442--451. \url{https://doi.org/10.1016/0005-2795(75)90109-9}
\bibitem{mckinney2010} McKinney, W. (2010). Data structures for statistical computing in Python. In S. van der Walt \& J. Millman (Eds.), \textit{Proceedings of the 9th Python in Science Conference} (pp. 56--61). \url{https://doi.org/10.25080/Majora-92bf1922-00a}
\bibitem{morales2025} Morales-Castro, J. A., y Espinosa Jiménez, P. M. (2025). Factores que afectan la sostenibilidad financiera de las Sociedades Financieras Populares de México. \textit{Entreciencias: Diálogos en la Sociedad del Conocimiento, 13}(27), 1--17. \url{https://doi.org/10.22201/enesl.20078064e.2025.27.89936}
\bibitem{morduch1999} Morduch, J. (1999). The microfinance promise. \textit{Journal of Economic Literature, 37}(4), 1569--1614. \url{https://doi.org/10.1257/jel.37.4.1569}
\bibitem{mullainathan2017} Mullainathan, S., \& Spiess, J. (2017). Machine learning: An applied econometric approach. \textit{Journal of Economic Perspectives, 31}(2), 87--106. \url{https://doi.org/10.1257/jep.31.2.87}
\bibitem{niculescu2005} Niculescu-Mizil, A., \& Caruana, R. (2005). Predicting good probabilities with supervised learning. In \textit{Proceedings of the 22nd International Conference on Machine Learning} (pp. 625--632). ACM. \url{https://doi.org/10.1145/1102351.1102430}
\bibitem{pedregosa2011} Pedregosa, F., Varoquaux, G., Gramfort, A., Michel, V., Thirion, B., Grisel, O., Blondel, M., Prettenhofer, P., Weiss, R., Dubourg, V., Vanderplas, J., Passos, A., Cournapeau, D., Brucher, M., Perrot, M., \& Duchesnay, É. (2011). Scikit-learn: Machine learning in Python. \textit{Journal of Machine Learning Research, 12}, 2825--2830. \url{https://jmlr.org/papers/v12/pedregosa11a.html}
\bibitem{prokhorenkova2018} Prokhorenkova, L., Gusev, G., Vorobev, A., Dorogush, A. V., \& Gulin, A. (2018). CatBoost: Unbiased boosting with categorical features. \textit{arXiv}. \url{https://doi.org/10.48550/arXiv.1706.09516}
\bibitem{raasveldt2019} Raasveldt, M., \& Mühleisen, H. (2019). DuckDB: An embeddable analytical database. In \textit{Proceedings of the 2019 International Conference on Management of Data} (pp. 1981--1984). ACM. \url{https://doi.org/10.1145/3299869.3320212}
\bibitem{robbins1951} Robbins, H., \& Monro, S. (1951). A stochastic approximation method. \textit{The Annals of Mathematical Statistics, 22}(3), 400--407. \url{https://doi.org/10.1214/aoms/1177729586}
\bibitem{ruder2016} Ruder, S. (2016). An overview of gradient descent optimization algorithms. \textit{arXiv}. \url{https://doi.org/10.48550/arXiv.1609.04747}
\bibitem{rudin2019} Rudin, C. (2019). Stop explaining black box machine learning models for high stakes decisions and use interpretable models instead. \textit{Nature Machine Intelligence, 1}(5), 206--215. \url{https://doi.org/10.1038/s42256-019-0048-x}
\bibitem{saito2015} Saito, T., \& Rehmsmeier, M. (2015). The precision-recall plot is more informative than the ROC plot when evaluating binary classifiers on imbalanced datasets. \textit{PLOS ONE, 10}(3), Article e0118432. \url{https://doi.org/10.1371/journal.pone.0118432}
\bibitem{scheffer2009} Scheffer, M., Bascompte, J., Brock, W. A., Brovkin, V., Carpenter, S. R., Dakos, V., Held, H., van Nes, E. H., Rietkerk, M., \& Sugihara, G. (2009). Early-warning signals for critical transitions. \textit{Nature, 461}(7260), 53--59. \url{https://doi.org/10.1038/nature08227}
\bibitem{seabold2010} Seabold, S., \& Perktold, J. (2010). Statsmodels: Econometric and statistical modeling with Python. In S. van der Walt \& J. Millman (Eds.), \textit{Proceedings of the 9th Python in Science Conference} (pp. 92--96). \url{https://doi.org/10.25080/Majora-92bf1922-011}
\bibitem{shapley1953} Shapley, L. S. (1953). A value for n-person games. In H. W. Kuhn \& A. W. Tucker (Eds.), \textit{Contributions to the theory of games} (Vol. 2, pp. 307--317). Princeton University Press.% VERIFICAR: datos de edición (no registrado en CrossRef).
\bibitem{shwartzziv2022} Shwartz-Ziv, R., \& Armon, A. (2022). Tabular data: Deep learning is not all you need. \textit{Information Fusion, 81}, 84--90. \url{https://doi.org/10.1016/j.inffus.2021.11.011}
\bibitem{simonin2020} Simonin, I., Brooks, M., \& Nieto Barajas, L. E. (2020). Portfolio recommendations to improve risk of default in microfinance. \textit{CIENCIA ergo-sum, 28}(1). \url{https://doi.org/10.30878/ces.v28n1a6}% VERIFICAR: número de artículo o páginas no listados en la página de la revista.
\bibitem{snoek2012} Snoek, J., Larochelle, H., \& Adams, R. P. (2012). Practical Bayesian optimization of machine learning algorithms. \textit{arXiv}. \url{https://doi.org/10.48550/arXiv.1206.2944}
\bibitem{stiglitz1981} Stiglitz, J. E., \& Weiss, A. (1981). Credit rationing in markets with imperfect information. \textit{The American Economic Review, 71}(3), 393--410.% VERIFICAR: DOI no registrado en CrossRef.
\bibitem{strumbelj2014} Štrumbelj, E., \& Kononenko, I. (2014). Explaining prediction models and individual predictions with feature contributions. \textit{Knowledge and Information Systems, 41}(3), 647--665. \url{https://doi.org/10.1007/s10115-013-0679-x}% Año de la edición impresa; CrossRef registra 2013 como publicación en línea.
\bibitem{tashman2000} Tashman, L. J. (2000). Out-of-sample tests of forecasting accuracy: An analysis and review. \textit{International Journal of Forecasting, 16}(4), 437--450. \url{https://doi.org/10.1016/S0169-2070(00)00065-0}
\bibitem{tibshirani1996} Tibshirani, R. (1996). Regression shrinkage and selection via the lasso. \textit{Journal of the Royal Statistical Society: Series B (Methodological), 58}(1), 267--288. \url{https://doi.org/10.1111/j.2517-6161.1996.tb02080.x}
\bibitem{torres2019} Torres Grimaldo, J. A., y Cano Morales, A. M. (2019). Importancia del gobierno corporativo en las sociedades financieras populares en México. \textit{Cuadernos de Contabilidad, 20}(49), 1--16. \url{https://doi.org/10.11144/javeriana.cc20-49.igcs}
\bibitem{vapnik2000} Vapnik, V. N. (2000). \textit{The nature of statistical learning theory} (2nd ed.). Springer. \url{https://doi.org/10.1007/978-1-4757-3264-1}
\bibitem{varian2014} Varian, H. R. (2014). Big data: New tricks for econometrics. \textit{Journal of Economic Perspectives, 28}(2), 3--28. \url{https://doi.org/10.1257/jep.28.2.3}
\bibitem{virtanen2020} Virtanen, P., Gommers, R., Oliphant, T. E., Haberland, M., Reddy, T., Cournapeau, D., Burovski, E., Peterson, P., Weckesser, W., Bright, J., van der Walt, S. J., Brett, M., Wilson, J., Millman, K. J., Mayorov, N., Nelson, A. R. J., Jones, E., Kern, R., Larson, E., \ldots{} Vázquez-Baeza, Y. (2020). SciPy 1.0: Fundamental algorithms for scientific computing in Python. \textit{Nature Methods, 17}(3), 261--272. \url{https://doi.org/10.1038/s41592-019-0686-2}
\bibitem{wilson1972} Wilson, D. L. (1972). Asymptotic properties of nearest neighbor rules using edited data. \textit{IEEE Transactions on Systems, Man, and Cybernetics, SMC-2}(3), 408--421. \url{https://doi.org/10.1109/TSMC.1972.4309137}
\bibitem{wolpert1992} Wolpert, D. H. (1992). Stacked generalization. \textit{Neural Networks, 5}(2), 241--259. \url{https://doi.org/10.1016/S0893-6080(05)80023-1}
\bibitem{wolpert1996} Wolpert, D. H. (1996). The lack of a priori distinctions between learning algorithms. \textit{Neural Computation, 8}(7), 1341--1390. \url{https://doi.org/10.1162/neco.1996.8.7.1341}
\bibitem{wooldridge2010} Wooldridge, J. M. (2010). \textit{Econometric analysis of cross section and panel data} (2nd ed.). MIT Press.% VERIFICAR: datos de edición del libro (no registrado en CrossRef).
\end{thebibliography}

% NOTA sobre C1: Cambiar "PRESENTA" por "PRESENTAN" en PORTADA.pdf (dos autores).
% Este cambio debe realizarse directamente en el archivo PDF original con un editor de PDF.

